\chapter{String de búsqueda y criterios de selección}
\label{anexo:a}

\section{Bases de datos consultadas}

\begin{itemize}
  \item IEEE Xplore Digital Library
  \item ACM Digital Library
  \item Scopus / ScienceDirect
  \item Google Scholar (complemento y trazabilidad de citas)
\end{itemize}

\section{Período y idioma}

\begin{itemize}
  \item \textbf{Período:} 2005--2026 (énfasis 2015--2026 para OpenAPI~3.x y BIAN Service Landscape).
  \item \textbf{Idioma:} principalmente inglés; se incluyen estudios en español cuando aportan al
        mapeo OpenAPI (p. ej. Casas et al., 2021).
\end{itemize}

\section{Cadenas de búsqueda (inglés)}

\begin{description}
  \item[Eje OpenAPI / APIs REST:]
    \begin{quote}\small\ttfamily
      (``OpenAPI'' OR ``Swagger'' OR ``OpenAPI Specification'') AND (``REST'' OR ``RESTful'')
      AND (``API'' OR ``web service'')
    \end{quote}
  \item[Eje alineación / matching:]
    \begin{quote}\small\ttfamily
      (``schema matching'' OR ``ontology alignment'' OR ``semantic similarity'' OR
      ``structural alignment'') AND (``OpenAPI'' OR ``API specification'')
    \end{quote}
  \item[Eje banca / arquitectura:]
    \begin{quote}\small\ttfamily
      (``BIAN'' OR ``banking architecture'' OR ``financial services'' OR ``open banking'') AND
      (``API'' OR ``service domain'' OR ``enterprise architecture'')
    \end{quote}
\end{description}

\section{Criterios de inclusión}

\begin{enumerate}
  \item Artículos de revisión sistemática, mapeo sistemático, \textit{survey} o revisión
        narrativa con método explícito.
  \item Estudios que aborden OpenAPI/\textit{Swagger}, matching estructural/semántico, ontologías,
        arquitectura bancaria o calidad de APIs REST.
  \item Publicaciones con acceso al texto completo (PDF institucional o repositorio del autor).
  \item Relevancia directa para delimitar el vacío OpenAPI--BIAN o para sustentar S, E o C.
\end{enumerate}

\section{Criterios de exclusión}

\begin{enumerate}
  \item Trabajos sin revisión por pares (posters, \textit{preprints} no validados) salvo
        referencia contextual.
  \item Estudios centrados exclusivamente en \textit{runtime}, latencia de \textit{gateway} o
        despliegue productivo sin aporte al diseño documental (delimitación E5).
  \item Duplicados (versiones de conferencia y journal del mismo estudio: se conserva la más
        completa).
  \item Papers cuya técnica no pueda vincularse al objeto de estudio (contrato OpenAPI bancario).
\end{enumerate}

\section{Resultado de la búsqueda}

La selección final para el estado del arte comprende \textbf{14 papers} agrupados en seis ejes
(E1--E6), documentados en el Anexo~B y sintetizados en el Capítulo~3 §3.5--§3.6. Los papers de
\textit{aporte} (22 en \texttt{bibliografia/aporte/}) se reservan para sustentar técnicas del
Capítulo~4 y no sustituyen al estado del arte.
